\section{Edges in a vertex set}
\label{sec:edgesIn}

Throughout, let $G$ be a simple graph on a vertex type $V$ and let
$S, T \subseteq V$. The combinatorial arguments that follow rest on a
single construction --- the set of edges of $G$ with both endpoints in
a vertex set $S$ --- and this section collects its basic properties:
monotonicity and the corner cases, the effect of deleting a vertex, and
the way the construction interacts with intersection and union.

\subsection{Definition and basic properties}

\begin{definition}[Edges induced on a vertex set]
  \label{def:edgesIn}
  \lean{SimpleGraph.edgesIn}
  \leanok
  For $S \subseteq V$, let
  \[
    \edgesIn{S} \;:=\; \{\,e \in E(G) : e \subseteq S\,\}
  \]
  be the set of edges of $G$ with both endpoints in $S$, that is, the
  edge set of the subgraph of $G$ induced on $S$. We say that $S$
  \emph{spans} these edges.
\end{definition}

\begin{lemma}[Monotonicity]
  \label{lem:edgesIn-mono}
  \lean{SimpleGraph.edgesIn_mono}
  \leanok
  \uses{def:edgesIn}
  Let $S \subseteq T \subseteq V$. Then
  $\edgesIn{S} \subseteq \edgesIn{T}$.
\end{lemma}
\begin{proof}
  \leanok
  An edge $e \subseteq S$ also satisfies $e \subseteq T$.
\end{proof}

\begin{lemma}[Corner cases]
  \label{lem:edgesIn-corners}
  \lean{SimpleGraph.edgesIn_univ, SimpleGraph.edgesIn_empty, SimpleGraph.edgesIn_singleton, SimpleGraph.edgesIn_bot}
  \leanok
  \uses{def:edgesIn}
  Let $v \in V$ and $S \subseteq V$. Then $\edgesIn{V} = E(G)$,
  $\edgesIn{\emptyset} = \emptyset$ and $\edgesIn{\{v\}} = \emptyset$.
  If $G$ has no edges, then $\edgesIn{S} = \emptyset$.
\end{lemma}
\begin{proof}
  \leanok
  All but the third are immediate from the definition. A singleton
  $\{v\}$ contains no edge, since $G$ has no loops and so the two
  endpoints of an edge are distinct.
\end{proof}

\subsection{Vertex deletion}

Let $v \in V$, and write $\{v\}^c = V \setminus \{v\}$. The set
$\edgesIn{\{v\}^c}$ consists of the edges of $G$ that remain when $v$
is deleted. Let $I_G(v) := \{\,e \in E(G) : v \in e\,\}$ be the set of
edges of $G$ incident to $v$. When $V$ is finite, $|I_G(v)|$ is the
degree of $v$. The next two lemmas show that deleting $v$ removes
exactly the edges of $I_G(v)$, and count what remains. The count is
used in \cref{lem:isLaman-two-le-degree} to show that every vertex of a
Laman graph on at least three vertices has degree at least two.

\begin{lemma}[Edges avoiding a vertex]
  \label{lem:mem-edgesIn-compl-singleton}
  \lean{SimpleGraph.mem_edgesIn_compl_singleton, SimpleGraph.edgesIn_compl_singleton}
  \leanok
  \uses{def:edgesIn}
  Let $v \in V$. Then
  \[
    \edgesIn{\{v\}^{c}} \;=\; E(G) \,\setminus\, I_G(v);
  \]
  that is, a pair $e$ of vertices lies in $\edgesIn{\{v\}^c}$ if and
  only if $e \in E(G)$ and $v \notin e$.
\end{lemma}
\begin{proof}
  \leanok
  An edge has both endpoints in $\{v\}^c$ exactly when neither
  endpoint is $v$, that is, when $v \notin e$. So
  $\edgesIn{\{v\}^c}$ consists of the edges of $G$ not incident to $v$,
  which is $E(G) \setminus I_G(v)$.
\end{proof}

\begin{lemma}[Vertex-deletion partition of edges]
  \label{lem:ncard-edgesIn-compl-singleton}
  \lean{SimpleGraph.ncard_edgesIn_compl_singleton_add_ncard_incidenceSet}
  \leanok
  \uses{lem:mem-edgesIn-compl-singleton}
  Let $V$ be finite and let $v \in V$. Then
  \[
    \bigl|\edgesIn{\{v\}^{c}}\bigr| \;+\; |I_G(v)| \;=\; |E(G)|.
  \]
\end{lemma}
\begin{proof}
  \leanok
  By \cref{lem:mem-edgesIn-compl-singleton}, $E(G)$ is the disjoint
  union of $\edgesIn{\{v\}^c}$ and $I_G(v)$, so the cardinalities add.
\end{proof}

\subsection{Modularity: intersection, union, and supermodular counts}

The map $S \mapsto \edgesIn{S}$ commutes with intersection but not with
union. An edge with one endpoint in $S \setminus T$ and the other in
$T \setminus S$ lies in $\edgesIn{S \cup T}$ but in neither
$\edgesIn{S}$ nor $\edgesIn{T}$. Counting edges therefore gives an
inequality (\cref{lem:edgesIn-ncard-add-le}). It becomes an identity,
Jord\'an's Lemma~2.1.1~\cite{jordan2016}, when the number of these
crossing edges is added to its left side. We need only the inequality,
which is the counting step in the proof that tight sets are closed
under union and intersection (\cref{thm:isTightOn-union-inter}).

\begin{lemma}[Distribution over intersection]
  \label{lem:edgesIn-inter}
  \lean{SimpleGraph.edgesIn_inter}
  \leanok
  \uses{def:edgesIn}
  Let $S, T \subseteq V$. Then
  $\edgesIn{S \cap T} = \edgesIn{S} \cap \edgesIn{T}$.
\end{lemma}
\begin{proof}
  \leanok
  An edge has both endpoints in $S \cap T$ if and only if both
  endpoints lie in $S$ and both lie in $T$.
\end{proof}

\begin{lemma}[Supermodularity of edge counts]
  \label{lem:edgesIn-ncard-add-le}
  \lean{SimpleGraph.edgesIn_ncard_add_le}
  \leanok
  \uses{lem:edgesIn-mono, lem:edgesIn-inter}
  Let $V$ be finite and let $S, T \subseteq V$. Then
  \[
    \bigl|\edgesIn{S}\bigr| + \bigl|\edgesIn{T}\bigr|
    \;\le\;
    \bigl|\edgesIn{S \cup T}\bigr| + \bigl|\edgesIn{S \cap T}\bigr|.
  \]
\end{lemma}
\begin{proof}
  \leanok
  Let $A = \edgesIn{S}$ and $B = \edgesIn{T}$, so that
  $|A| + |B| = |A \cup B| + |A \cap B|$. By \cref{lem:edgesIn-inter},
  $A \cap B = \edgesIn{S \cap T}$. By monotonicity
  (\cref{lem:edgesIn-mono}), $A \cup B \subseteq \edgesIn{S \cup T}$, so
  $|A \cup B| \le |\edgesIn{S \cup T}|$. Together these give the
  inequality.
\end{proof}

\section{\texorpdfstring{$(k,\ell)$}{(k,l)}-sparsity and \texorpdfstring{$(k,\ell)$}{(k,l)}-tightness}
\label{sec:sparsity}

This section introduces $(k,\ell)$-sparsity and $(k,\ell)$-tightness and
develops their basic theory: the global edge bound and its low-degree
consequences, invariance under graph isomorphism, and the fact that the
tight subsets of a sparse graph are closed under union and intersection.
A final subsection isolates the sparse-graph counting fact behind the
reverse Type~II Henneberg move. Throughout this section $k$ and $\ell$
are non-negative integers (think $k = 2$, $\ell = 3$ for the Laman case
in \cref{sec:intro}).

\subsection{Definitions}

\begin{definition}[$(k,\ell)$-sparse]
  \label{def:isSparse}
  \lean{SimpleGraph.IsSparse}
  \leanok
  \uses{def:edgesIn}
  $G$ is \emph{$(k,\ell)$-sparse} if every finite set $s \subseteq V$
  with $\ell \le k\,|s|$ spans at most $k\,|s| - \ell$ edges, that is,
  \[
    \bigl|\edgesIn{s}\bigr| + \ell \;\le\; k\,|s|.
  \]
\end{definition}
\begin{fmlnote}
  The formal definition states the bound in the additive form displayed
  above. Edge counts there are natural numbers, in which
  $k\,|s| - \ell$ would be truncated at $0$. Under the hypothesis
  $\ell \le k\,|s|$ the two forms agree.
\end{fmlnote}

\begin{definition}[$(k,\ell)$-tight]
  \label{def:isTight}
  \lean{SimpleGraph.IsTight}
  \leanok
  \uses{def:isSparse}
  Let $V$ be finite. The graph $G$ is \emph{$(k,\ell)$-tight} if it is
  $(k,\ell)$-sparse and
  \[
    |E(G)| + \ell \;=\; k\,|V|.
  \]
\end{definition}

\begin{definition}[Tight subset]
  \label{def:isTightOn}
  \lean{SimpleGraph.IsTightOn}
  \leanok
  \uses{def:edgesIn}
  A finite set $s \subseteq V$ is \emph{$(k,\ell)$-tight in $G$} if it
  attains the sparsity bound:
  \[
    \bigl|\edgesIn{s}\bigr| + \ell \;=\; k\,|s|.
  \]
\end{definition}

Tightness of a set is the local form of tightness of the graph. Since
$\edgesIn{V} = E(G)$ (\cref{lem:edgesIn-corners}), a $(k,\ell)$-sparse
graph $G$ on a finite set $V$ is $(k,\ell)$-tight exactly when $V$ is a
$(k,\ell)$-tight set in $G$. At $(k,\ell) = (2,3)$ the hypothesis
$\ell \le k\,|s|$ of \cref{def:isSparse} says that $|s| \ge 2$. So a
finite simple graph is $(2,3)$-sparse exactly when it is sparse in the
sense of Jord\'an~\cite{jordan2016}~\S 2.1, and its $(2,3)$-tight sets
are what he calls \emph{critical sets}. We use ``tight'' throughout.

\subsection{Basic properties and the global edge bound}

Sparsity passes to subgraphs. Applied to the whole vertex set, it
bounds the number of edges of $G$, so that a tight graph has as many
edges as a sparse graph on $V$ can have.

\begin{lemma}[Subgraphs of sparse graphs are sparse]
  \label{lem:isSparse-mono-left}
  \lean{SimpleGraph.IsSparse.mono_left}
  \leanok
  \uses{def:isSparse}
  Let $G'$ be a graph on $V$ with $E(G) \subseteq E(G')$. If $G'$ is
  $(k,\ell)$-sparse, then so is $G$.
\end{lemma}
\begin{proof}
  \leanok
  For a finite set $s \subseteq V$, every edge of $G$ with both
  endpoints in $s$ is an edge of $G'$ with both endpoints in $s$, so
  $|\edgesIn[G]{s}| \le |\edgesIn[G']{s}|$. The sparsity bound for
  $G'$ at $s$ therefore bounds $|\edgesIn[G]{s}|$.
\end{proof}

\begin{lemma}[Global edge bound]
  \label{lem:isSparse-edgeSet-bound}
  \lean{SimpleGraph.IsSparse.edgeSet_ncard_add_le}
  \leanok
  \uses{def:isSparse}
  Let $V$ be finite, and let $G$ be $(k,\ell)$-sparse with
  $\ell \le k\,|V|$. Then
  \[
    |E(G)| + \ell \;\le\; k\,|V|.
  \]
\end{lemma}
\begin{proof}
  \leanok
  Apply the sparsity bound at $s = V$ and use $\edgesIn{V} = E(G)$
  (\cref{lem:edgesIn-corners}).
\end{proof}

\begin{lemma}[A proper supergraph of a tight graph is not sparse]
  \label{lem:isTight-not-isSparse-of-lt}
  \lean{SimpleGraph.IsTight.not_isSparse_of_lt}
  \leanok
  \uses{def:isTight, lem:isSparse-edgeSet-bound}
  Let $V$ be finite, let $G$ be $(k,\ell)$-tight, and let $H$ be a
  graph on $V$ with $E(G) \subsetneq E(H)$. Then $H$ is not
  $(k,\ell)$-sparse.
\end{lemma}
\begin{proof}
  \leanok
  By tightness $|E(G)| + \ell = k\,|V|$; in particular
  $\ell \le k\,|V|$. Since $E(G)$ is a proper subset of $E(H)$, we have
  $|E(G)| < |E(H)|$, and so $|E(H)| + \ell > k\,|V|$. By
  \cref{lem:isSparse-edgeSet-bound}, $H$ is not $(k,\ell)$-sparse.
\end{proof}

\subsection{Low-degree vertex and non-adjacency consequences}

Both lemmas below concern $(2,3)$-sparse graphs. The first is applied
to Laman graphs in \cref{lem:isLaman-exists-two-le-degree-le-three}.
The second is used in the reverse decomposition of a $(2,3)$-sparse
graph, at a vertex of degree three
(\cref{thm:isSparse-exists-typeI-or-typeII-reverse}).

\begin{lemma}[A vertex of degree at most 3]
  \label{lem:isSparse-exists-degree-le-three}
  \lean{SimpleGraph.IsSparse.exists_degree_le_three}
  \leanok
  \uses{def:isSparse, lem:isSparse-edgeSet-bound}
  Let $V$ be finite with $|V| \ge 2$, and let $G$ be $(2,3)$-sparse.
  Then some vertex of $G$ has degree at most $3$.
\end{lemma}
\begin{proof}
  \leanok
  Suppose that every vertex has degree at least $4$. By the handshake
  identity, $2|E(G)| = \sum_v \deg_G(v) \ge 4|V|$. The global edge
  bound (\cref{lem:isSparse-edgeSet-bound}) applies at
  $(k, \ell) = (2, 3)$, since $3 \le 2|V|$ when $|V| \ge 2$. It gives
  $|E(G)| + 3 \le 2|V|$, so $2|E(G)| \le 4|V| - 6 < 4|V|$, a
  contradiction.
\end{proof}

\begin{lemma}[Non-adjacent pair among three neighbors]
  \label{lem:isSparse-exists-nonadj-among-three-neighbors}
  \lean{SimpleGraph.IsSparse.exists_nonadj_among_three_neighbors}
  \leanok
  \uses{def:isSparse}
  Let $V$ be finite and let $G$ be $(2,3)$-sparse. Suppose that
  $v \in V$ has three pairwise distinct neighbors $a, b, c$. Then at
  least one of the pairs $\{a, b\}$, $\{a, c\}$, $\{b, c\}$ is
  non-adjacent in $G$.
\end{lemma}
\begin{proof}
  \leanok
  Suppose that all three pairs are adjacent, and let
  $X = \{v, a, b, c\}$. The vertices $a$, $b$ and $c$ are distinct, and
  each differs from $v$ because $G$ has no loops, so $|X| = 4$. The set
  $X$ spans the six distinct edges $va$, $vb$, $vc$, $ab$, $ac$ and
  $bc$, so $|\edgesIn{X}| \ge 6$. But $(2,3)$-sparsity applies at $X$,
  since $3 \le 2 \cdot 4$, and gives
  $|\edgesIn{X}| \le 2 \cdot 4 - 3 = 5$, a contradiction.
\end{proof}

\subsection{Transport along graph isomorphism}

Both sparsity and tightness are isomorphism invariants. We first record
how $\edgesIn{\cdot}$ transports along an isomorphism, then deduce the
invariance of each predicate.

\begin{lemma}[$\edgesIn{\cdot}$ transports along isomorphisms]
  \label{lem:iso-image-edgesIn}
  \lean{SimpleGraph.Iso.image_edgesIn}
  \leanok
  \uses{def:edgesIn}
  Let $\varphi : G \simeq H$ be a graph isomorphism and $s \subseteq V$,
  and write $\varphi$ also for the induced map on edges (unordered pairs
  of vertices). Then
  \[
    \varphi\bigl[\edgesIn[G]{s}\bigr] = \edgesIn[H]{\varphi(s)}.
  \]
\end{lemma}
\begin{proof}
  \leanok
  An unordered pair $\{u, v\}$ in the image set comes from a pair
  $\{p, q\} \subseteq s$ in $\edgesIn[G]{s}$ with $u = \varphi(p)$ and
  $v = \varphi(q)$, by checking the pair's two representatives;
  $\varphi$ preserves adjacency in each direction.
\end{proof}

\begin{lemma}[Sparsity is preserved by isomorphism]
  \label{lem:isSparse-iso}
  \lean{SimpleGraph.IsSparse.iso}
  \leanok
  \uses{def:isSparse, lem:iso-image-edgesIn}
  If $\varphi : G \simeq H$ and $G$ is $(k,\ell)$-sparse, then so is $H$.
\end{lemma}
\begin{proof}
  \leanok
  For $s \subseteq W$ a finite set on the codomain, transport along
  $\varphi^{-1}$ to a set $s' \subseteq V$ of the same cardinality.
  By \cref{lem:iso-image-edgesIn},
  $|\edgesIn[H]{s}| = |\edgesIn[G]{s'}|$. Apply $(k,\ell)$-sparsity of
  $G$ at $s'$.
\end{proof}

\begin{lemma}[Tightness is preserved by isomorphism]
  \label{lem:isTight-iso}
  \lean{SimpleGraph.IsTight.iso}
  \leanok
  \uses{def:isTight, lem:isSparse-iso}
  If $\varphi : G \simeq H$ and $G$ is $(k,\ell)$-tight, then so is $H$.
\end{lemma}
\begin{proof}
  \leanok
  Sparsity transports by \cref{lem:isSparse-iso}. For the global edge
  identity, $\varphi$ induces bijections $E(G) \cong E(H)$ and
  $V \cong W$, so $|E(H)| = |E(G)|$ and $|W| = |V|$. Substituting these
  into the tightness equation for $G$ yields the one for $H$.
\end{proof}

\subsection{The tight-subset lattice}

The squeeze lemma below is a converse to the sparsity inequality
within a single sparse graph: a vertex set whose edge count reaches the
sparsity bound from below must attain it with equality, hence is tight.
It is the edge-counting input used in
\cref{lem:isSparse-typeII-reverse-blocker} to extract a tight set from
a sparsity violation on a candidate graph.

\begin{lemma}[Squeeze]
  \label{lem:isSparse-isTightOn-of-le}
  \lean{SimpleGraph.IsSparse.isTightOn_of_le}
  \leanok
  \uses{def:isSparse, def:isTightOn}
  Let $G$ be $(k,\ell)$-sparse, $s \subseteq V$ finite with
  $\ell \le k\,|s|$, and suppose
  \[
    k\,|s| \;\le\; \bigl|\edgesIn{s}\bigr| + \ell.
  \]
  Then $s$ is $(k,\ell)$-tight in $G$.
\end{lemma}
\begin{proof}
  \leanok
  The sparsity bound at $s$ gives the reverse inequality
  $|\edgesIn{s}| + \ell \le k\,|s|$. The two together force equality.
\end{proof}

Inside a fixed $(k,\ell)$-sparse graph, the family of tight subsets is
closed under union and intersection, subject to a size proviso on the
intersection.

\begin{theorem}[Tight subsets are closed under union and intersection]
  \label{thm:isTightOn-union-inter}
  \lean{SimpleGraph.IsTightOn.union_inter}
  \leanok
  \uses{def:isSparse, def:isTightOn, lem:edgesIn-ncard-add-le}
  Let $V$ be finite and $G$ be $(k,\ell)$-sparse. Let $s, t \subseteq V$
  be finite vertex sets, both $(k,\ell)$-tight in $G$, and suppose
  $\ell \le k\,|s \cap t|$. Then both $s \cup t$ and $s \cap t$ are
  $(k,\ell)$-tight in $G$.
\end{theorem}
\begin{proof}
  \leanok
  Sparsity applied to $s \cap t$ (using the size hypothesis) and to
  $s \cup t$ (whose cardinality is at least $|s \cap t|$) gives
  \begin{align*}
    \bigl|\edgesIn{s \cap t}\bigr| + \ell &\le k\,|s \cap t|, \\
    \bigl|\edgesIn{s \cup t}\bigr| + \ell &\le k\,|s \cup t|.
  \end{align*}
  Adding these and using
  $|s \cup t| + |s \cap t| = |s| + |t|$
  yields
  \[
    \bigl|\edgesIn{s \cup t}\bigr|
      + \bigl|\edgesIn{s \cap t}\bigr|
      + 2\ell
    \;\le\; k\,|s| + k\,|t|.
  \]

  On the other hand, supermodularity
  (\cref{lem:edgesIn-ncard-add-le}) gives
  \[
    \bigl|\edgesIn{s}\bigr| + \bigl|\edgesIn{t}\bigr|
    \;\le\;
    \bigl|\edgesIn{s \cup t}\bigr|
      + \bigl|\edgesIn{s \cap t}\bigr|,
  \]
  and tightness of $s$ and $t$ rewrites the left side as
  $k\,|s| + k\,|t| - 2\ell$. Chaining the two inequalities forces
  every inequality to be an equality, i.e.\ both $s \cup t$ and
  $s \cap t$ are tight.
\end{proof}

Two refinements of \cref{thm:isTightOn-union-inter}'s union half drop
the size proviso on $s \cap t$ in exchange for an explicit set $F$ of
extra edges that closes the accounting. They specialize to the union
half of \cref{thm:isTightOn-union-inter} at $F = \emptyset$.

\begin{theorem}[Tight union with bonus edges]
  \label{thm:isTightOn-union-with-bonus}
  \lean{SimpleGraph.IsTightOn.union_with_bonus}
  \leanok
  \uses{def:isSparse, def:isTightOn, lem:edgesIn-ncard-add-le}
  Let $V$ be finite and $G$ be $(k,\ell)$-sparse. Let
  $s_1, s_2 \subseteq V$ be both $(k,\ell)$-tight, and let
  $F \subseteq \edgesIn{s_1 \cup s_2}$ be a finite set of edges
  disjoint from $\edgesIn{s_1} \cup \edgesIn{s_2}$, satisfying
  \[
    |\edgesIn{s_1 \cap s_2}| + \ell \;\le\; |F| + k\,|s_1 \cap s_2|.
  \]
  Then $s_1 \cup s_2$ is $(k,\ell)$-tight.
\end{theorem}
\begin{proof}
  \leanok
  Refined supermodularity, applied to the disjoint inclusion
  $\edgesIn{s_1} \cup \edgesIn{s_2} \cup F \subseteq \edgesIn{s_1 \cup s_2}$
  together with $\edgesIn{s_1} \cap \edgesIn{s_2} = \edgesIn{s_1 \cap s_2}$,
  gives
  \[
    |\edgesIn{s_1}| + |\edgesIn{s_2}| + |F|
    \;\le\;
    |\edgesIn{s_1 \cup s_2}| + |\edgesIn{s_1 \cap s_2}|.
  \]
  Tightness of $s_1$ and $s_2$ rewrites the left side as
  $k\,|s_1| + k\,|s_2| - 2\ell + |F|$. Combined with the
  $|s_1| + |s_2| = |s_1 \cup s_2| + |s_1 \cap s_2|$ identity and the
  hypothesis on $F$, this lower-bounds $|\edgesIn{s_1 \cup s_2}|$ by
  $k\,|s_1 \cup s_2| - \ell$, which against sparsity forces equality.
\end{proof}

\begin{theorem}[Tight extension by a vertex with $k$ boundary edges]
  \label{thm:isTightOn-insert-vertex-with-edges}
  \lean{SimpleGraph.IsTightOn.insert_vertex_with_edges}
  \leanok
  \uses{def:isSparse, def:isTightOn}
  Let $V$ be finite and $G$ be $(k,\ell)$-sparse. Let
  $s \subseteq V$ be $(k,\ell)$-tight, $w \notin s$, and let
  $F \subseteq \edgesIn{s \cup \{w\}}$ be a finite set of edges
  disjoint from $\edgesIn{s}$ with $|F| \ge k$. Then $s \cup \{w\}$
  is $(k,\ell)$-tight.
\end{theorem}
\begin{proof}
  \leanok
  The disjoint inclusion
  $\edgesIn{s} \cup F \subseteq \edgesIn{s \cup \{w\}}$ gives
  $|\edgesIn{s \cup \{w\}}| \ge |\edgesIn{s}| + |F|$. Tightness of $s$
  and $|F| \ge k$ rewrite this as
  $|\edgesIn{s \cup \{w\}}| \ge k\,|s| - \ell + k = k\,|s \cup \{w\}| - \ell$,
  which against sparsity at $s \cup \{w\}$ forces equality.
\end{proof}

\subsection{Per-pair tight-blocker witness (Type~II reverse, sparse form)}
\label{sec:isSparse-typeII-reverse-blocker}

The reverse of the Type~II Henneberg move rests on the following
sparse-graph fact: a sparsity violation on the bridge-augmented
candidate graph $G' := G[\{w \ne v\}] \cup \{\{x, y\}\}$ yields a
$(2, 3)$-tight set in $G$ containing both $x$ and $y$. It is the
edge-counting input to the sparse-graph reverse decomposition
(\cref{thm:isSparse-exists-typeI-or-typeII-reverse}) and, through it,
the Laman reverse decomposition
(\cref{thm:isLaman-exists-typeI-or-typeII-reverse}).

\begin{lemma}[Per-pair tight blocker, sparse form]
  \label{lem:isSparse-typeII-reverse-blocker}
  \lean{SimpleGraph.IsSparse.typeII_reverse_blocker}
  \leanok
  \uses{def:isSparse, def:isTightOn,
    lem:isSparse-isTightOn-of-le}
  Let $V$ be finite, let $G$ be $(2, 3)$-sparse on $V$, and fix
  $v \in V$ and $x, y \in V$ both $\ne v$. Suppose the bridge-augmented
  candidate $G' := G[\{w \ne v\}] \cup \{\{x, y\}\}$ is \emph{not}
  $(2, 3)$-sparse. Then there exists a finite $S \subseteq V$ with
  $v \notin S$, $x, y \in S$, and $S$ $(2, 3)$-tight in $G$.
\end{lemma}
\begin{proof}
  \leanok
  From the failure of $(2,3)$-sparsity of $G'$, extract a finite
  violating set $s'$ of vertices of $G'$, and let
  $S \subseteq V \setminus \{v\}$ be its image under the inclusion of
  $G'$'s vertices into $V$. The edges of $G'$ on $s'$ correspond to the
  edges of $G$ on $S$, together with the bridge $\{x, y\}$ when both
  $x, y \in s'$.

  Split on whether $x, y \in s'$. In the bridge-included case the
  overcount is at most one edge, which by the squeeze lemma
  (\cref{lem:isSparse-isTightOn-of-le}) forces $S$ to be $(2,3)$-tight
  in $G$. In the bridge-excluded case the bridge contributes nothing,
  so the count already violates $(2,3)$-sparsity of $G$, contradicting
  the hypothesis; this case therefore cannot occur.
\end{proof}
\begin{fmlnote}
  The correspondence between the edges of $G'$ on $s'$ and the edges of
  $G$ on $S$ rests on a project-internal lemma about the induced
  subgraph $G[\{w \ne v\}]$: the image of its edge set on $s'$ under the
  vertex inclusion is $\edgesIn[G]{S}$, and the matching count uses that
  this edge map is injective, which follows from injectivity of the
  vertex inclusion.
\end{fmlnote}
